% Options for packages loaded elsewhere
\PassOptionsToPackage{unicode}{hyperref}
\PassOptionsToPackage{hyphens}{url}
\documentclass[
  a4paper,
]{ltjsarticle}
\usepackage{xcolor}
\usepackage[margin=25mm]{geometry}
\usepackage{amsmath,amssymb}
\setcounter{secnumdepth}{-\maxdimen} % remove section numbering
\usepackage{iftex}
\ifPDFTeX
  \usepackage[T1]{fontenc}
  \usepackage[utf8]{inputenc}
  \usepackage{textcomp} % provide euro and other symbols
\else % if luatex or xetex
  \usepackage{unicode-math} % this also loads fontspec
  \defaultfontfeatures{Scale=MatchLowercase}
  \defaultfontfeatures[\rmfamily]{Ligatures=TeX,Scale=1}
\fi
\usepackage{lmodern}
\ifPDFTeX\else
  % xetex/luatex font selection
\fi
% Use upquote if available, for straight quotes in verbatim environments
\IfFileExists{upquote.sty}{\usepackage{upquote}}{}
\IfFileExists{microtype.sty}{% use microtype if available
  \usepackage[]{microtype}
  \UseMicrotypeSet[protrusion]{basicmath} % disable protrusion for tt fonts
}{}
\makeatletter
\@ifundefined{KOMAClassName}{% if non-KOMA class
  \IfFileExists{parskip.sty}{%
    \usepackage{parskip}
  }{% else
    \setlength{\parindent}{0pt}
    \setlength{\parskip}{6pt plus 2pt minus 1pt}}
}{% if KOMA class
  \KOMAoptions{parskip=half}}
\makeatother
\usepackage{longtable,booktabs,array}
\newcounter{none} % for unnumbered tables
\usepackage{calc} % for calculating minipage widths
% Correct order of tables after \paragraph or \subparagraph
\usepackage{etoolbox}
\makeatletter
\patchcmd\longtable{\par}{\if@noskipsec\mbox{}\fi\par}{}{}
\makeatother
% Allow footnotes in longtable head/foot
\IfFileExists{footnotehyper.sty}{\usepackage{footnotehyper}}{\usepackage{footnote}}
\makesavenoteenv{longtable}
\setlength{\emergencystretch}{3em} % prevent overfull lines
\providecommand{\tightlist}{%
  \setlength{\itemsep}{0pt}\setlength{\parskip}{0pt}}
% Glyph map for lualatex (newunicodechar). Maps symbols that may lack font glyphs.
\usepackage{newunicodechar}
\usepackage{amssymb}
\newunicodechar{∎}{\ensuremath{\blacksquare}}
\newunicodechar{✓}{\ensuremath{\checkmark}}
\newunicodechar{✗}{\ensuremath{\times}}
\newunicodechar{⟹}{\ensuremath{\Longrightarrow}}
\newunicodechar{⟺}{\ensuremath{\Longleftrightarrow}}
\newunicodechar{↪}{\ensuremath{\hookrightarrow}}
\newunicodechar{⌊}{\ensuremath{\lfloor}}
\newunicodechar{⌋}{\ensuremath{\rfloor}}
\newunicodechar{≃}{\ensuremath{\simeq}}
\newunicodechar{≈}{\ensuremath{\approx}}
\newunicodechar{≤}{\ensuremath{\leq}}
\newunicodechar{≥}{\ensuremath{\geq}}
\newunicodechar{≠}{\ensuremath{\neq}}
\newunicodechar{→}{\ensuremath{\rightarrow}}
\newunicodechar{←}{\ensuremath{\leftarrow}}
\newunicodechar{↔}{\ensuremath{\leftrightarrow}}
\newunicodechar{⊥}{\ensuremath{\perp}}
\newunicodechar{√}{\ensuremath{\surd}}
\newunicodechar{∑}{\ensuremath{\sum}}
\newunicodechar{∏}{\ensuremath{\prod}}
\newunicodechar{∞}{\ensuremath{\infty}}
\newunicodechar{∈}{\ensuremath{\in}}
\newunicodechar{∀}{\ensuremath{\forall}}
\newunicodechar{∣}{\ensuremath{\mid}}
\newunicodechar{γ}{\ensuremath{\gamma}}
\newunicodechar{λ}{\ensuremath{\lambda}}
\newunicodechar{μ}{\ensuremath{\mu}}
\newunicodechar{ρ}{\ensuremath{\rho}}
\newunicodechar{σ}{\ensuremath{\sigma}}
\newunicodechar{φ}{\ensuremath{\varphi}}
\newunicodechar{ψ}{\ensuremath{\psi}}
\newunicodechar{θ}{\ensuremath{\theta}}
\newunicodechar{Δ}{\ensuremath{\Delta}}
\newunicodechar{Π}{\ensuremath{\Pi}}
\newunicodechar{Λ}{\ensuremath{\Lambda}}
\newunicodechar{τ}{\ensuremath{\tau}}
\newunicodechar{ε}{\ensuremath{\varepsilon}}
\newunicodechar{ω}{\ensuremath{\omega}}
\newunicodechar{π}{\ensuremath{\pi}}
\newunicodechar{ν}{\ensuremath{\nu}}
\newunicodechar{±}{\ensuremath{\pm}}
\newunicodechar{×}{\ensuremath{\times}}
\newunicodechar{·}{\ensuremath{\cdot}}
\usepackage{bookmark}
\IfFileExists{xurl.sty}{\usepackage{xurl}}{} % add URL line breaks if available
\urlstyle{same}
\hypersetup{
  pdftitle={Zero-Closure Factorization of High-Symmetry Floors --- Complete \textbackslash pm i Photon-Pair Covering of Even Systems N=8k and Minimal p-Wave Condensates of Odd Systems},
  pdfauthor={Noriaki Kihara (WF System Co., Ltd.), ORCID 0009-0004-6753-4020},
  hidelinks,
  pdfcreator={LaTeX via pandoc}}

\title{Zero-Closure Factorization of High-Symmetry Floors --- Complete
\(\pm i\) Photon-Pair Covering of Even Systems \(N=8k\) and Minimal
\(p\)-Wave Condensates of Odd Systems}
\author{Noriaki Kihara (WF System Co., Ltd.), ORCID 0009-0004-6753-4020}
\date{2026-09-12}

\begin{document}
\maketitle

\textbf{Series:} Foundations of Self-Consistent Relational-Wave Closed
Systems (Paper 4 of 10)\\
\textbf{Author:} Noriaki Kihara (WF System Co., Ltd.)　\textbf{Date:}
2026-09-12\\
\textbf{Version DOI:} 10.5281/zenodo.22729081\\
\textbf{Concept DOI:} 10.5281/zenodo.22729080\\
\textbf{ORCID:} 0009-0004-6753-4020\\
\textbf{Zenodo:} https://zenodo.org/records/22729081\\

\begin{center}\rule{0.5\linewidth}{0.5pt}\end{center}

\subsection{Abstract}\label{abstract}

We study the square-closure \(\sum_e z_e^2=0\) of the high-symmetry
floor as a factorization into smaller zero-closure blocks, and unify it
into a \textbf{factorization theorem of cyclic square waves} in which
the minimal block size is determined by the number theory of \(N\). We
reconstruct the cyclic-Fourier relative-equilibrium floor
\(z_{ij}=r_d e^{i2\pi(i+j)/N}\) from its defining equation, and
exhaustively audit all unordered pairs, 3-waves, and 4-waves.
\textbf{Theorem:} within a fixed distance class, the square wave
\(w=z^2\) cycles through a regular \(L\)-gon with phase step
\(q=e^{i8\pi/N}\) (order \(L=N/\gcd(N,4)\)), and the minimal
zero-closure block is the regular \(p\)-gon with
\(p=\operatorname{spf}(L)\) (smallest prime factor of \(L\)). Hence
\(p=2\iff L\) even \(\iff 8\mid N\)------\textbf{an exact two-wave
self-closure \(z_b=\pm i\,z_a\) (photon-type \(\pm i\) pair) within one
and the same distance class is possible only for \(N=8k\) (proved for
arbitrary \(N\))}. In the full space (including across distinct distance
classes) the same selection rule was confirmed by an exhaustive pair
audit over all \(N=3\)--\(40\) (via the phase lattice + parity +
amplitude reflection \(r_d=r_{N/2-d}\); a general proof for arbitrary
\(N\) that there is no accidental equal-amplitude pair other than the
reflection is a future task, §8). For \(N=8k\), using
\(z_{i+m}=\pm i\,z_i\) with \(m=N/8\), all \(M\) waves can be
\textbf{explicitly and completely matched} into \(M/2\) photon pairs
without any graph search (closed form of the candidate edge count
\(E_\gamma=N(4N-11)/2=84,424,1020,1872,2980\); the candidate graph is
bipartite with degrees \(2/4/8\) and connected-component sizes
\(4/8/16\), and the matching is non-unique). Even \(N\) other than
\(N=8k\) can be \(100\%\) decomposed into minimal \(p\)-gon
zero-closures, and each distance class \(C_d\) itself independently
satisfies \(\sum_{e\in C_d}z_e^2=0\). Odd \(N\) have zero two-wave
pairs, and the minimal prime-factor \(p\)-wave (\(3\mid N\) gives
3-waves, odd primes \(N\) give \(N\)-waves) is the fundamental unit. The
factorization is in general non-unique. On the floor all of these blocks
are phase-locked to a single carrier with frequency ratio \(1{:}1\)
(algebraically decomposable \(\ne\) dynamically independent).
Classification: existence of a static algebraic factorization (particle
identification and dynamical conservation are separate tasks).

\begin{center}\rule{0.5\linewidth}{0.5pt}\end{center}

\subsection{1. The Problem --- Factorization of
Closure}\label{the-problem-factorization-of-closure}

The floors of Papers 1 and 3 satisfy \(\sum_e z_e^2=0\) (zero-closure).
We ask whether this total closure can be \textbf{factorized} into a
direct sum of smaller zero-closure blocks, and how the type of those
blocks varies with \(N\). The minimal nontrivial block is a 2-wave
(\(z_2=\pm i z_1\), photon-type \(\pm i\) pair) (Paper 9).

\subsection{2. Objects and Methods}\label{objects-and-methods}

We reconstruct the high-symmetry theoretical floor
\(z_{ij}=r_d e^{i2\pi(i+j)/N}\) (case (C) of Paper 2, cyclic-Fourier
relative equilibrium) from its defining equation. In the full space we
independently cross-check \(\|z\|^2\), \(\sum z\), \(\sum z^2\), and the
eigen-residual. The two-wave condition is exhaustively checked over all
pairs with the dimensionless residual
\(\varepsilon_{ij}=|z_i^2+z_j^2|/(|z_i|^2+|z_j|^2)\) (strict is
\(\varepsilon\le10^{-12}\)); 3-waves and 4-waves are exhaustively
checked by meet-in-the-middle. With \(Z_0=(v+\delta g)/\|\cdot\|\) and
\(\delta=10^{-15}\), the \(N\) values and counts under strict judgment
are unchanged.

\subsection{\texorpdfstring{3. Factorization Theorem of Cyclic Square
Waves
(\(L=N/\gcd(N,4)\))}{3. Factorization Theorem of Cyclic Square Waves (L=N/\textbackslash gcd(N,4))}}\label{factorization-theorem-of-cyclic-square-waves-lngcdn4}

Since the edge \(\{i,i+d\}\) of a fixed distance class \(C_d\) is
\(z_{i,i+d}=r_d e^{i2\pi(2i+d)/N}\), the square wave

\[w_i=z_{i,i+d}^2=r_d^2\,e^{i4\pi(2i+d)/N},\qquad w_{i+1}=q\,w_i,\quad q=e^{i8\pi/N},\quad L=\operatorname{ord}(q)=\frac{N}{\gcd(N,4)}.\]

Since the square waves within a class cycle through the vertices of a
regular \(L\)-gon, \(q^{L/p}\) spans a subgroup of order
\(p=\operatorname{spf}(L)\) = a regular \(p\)-gon, whose vertex sum
\(=0\)------\textbf{a \(p\)-wave zero-closure block exists}.
Furthermore, by the \textbf{vanishing-sum theorem of Lam--Leung
{[}2{]}}, the minimal weight (number of terms) of a nonempty vanishing
sum of \(L\)-th roots equals the smallest prime factor of \(L\). Since
the square waves within a distance class are \(L\)-th roots of equal
modulus \(r_d^2\), this \(p\)-gon is not merely a construction example
but \textbf{the minimal zero-closure block within the same class} (not
only existence but minimality). This makes a single formula explain both
even and odd cases:

{\def\LTcaptype{none} % do not increment counter
\begin{longtable}[]{@{}
  >{\raggedright\arraybackslash}p{(\linewidth - 4\tabcolsep) * \real{0.3333}}
  >{\raggedright\arraybackslash}p{(\linewidth - 4\tabcolsep) * \real{0.3333}}
  >{\raggedright\arraybackslash}p{(\linewidth - 4\tabcolsep) * \real{0.3333}}@{}}
\toprule\noalign{}
\begin{minipage}[b]{\linewidth}\raggedright
\(N\)
\end{minipage} & \begin{minipage}[b]{\linewidth}\raggedright
\(L=N/\gcd(N,4)\)
\end{minipage} & \begin{minipage}[b]{\linewidth}\raggedright
Minimal block \(p=\operatorname{spf}(L)\)
\end{minipage} \\
\midrule\noalign{}
\endhead
\bottomrule\noalign{}
\endlastfoot
odd & \(N\) & \(\operatorname{spf}(N)\) (\(3\mid N\)→3-wave,
\(N=25,35\)→5-wave, odd prime \(N\)→\(N\)-wave) \\
\(N\equiv2\ (\mathrm{mod}\,4)\) & \(N/2\) &
\(\operatorname{spf}(N/2)\) \\
\(N\equiv4\ (\mathrm{mod}\,8)\) & \(N/4\) (odd) &
\(\operatorname{spf}(N/4)\) \\
\(N\equiv0\ (\mathrm{mod}\,8)\) & \(N/4\) (even) & \(2\) (\(\pm i\)
photon pair) \\
\end{longtable}
}

Hence \(\boxed{p=2\iff L\ \text{even}\iff 8\mid N}\)------\textbf{a
two-wave zero-closure (\(\pm i\) photon pair) within the same distance
class is possible only for \(8\mid N\)}. \(\operatorname{ord}(q)=L\) and
the \(p\)-gon sum \(=0\) were confirmed for all \(N=3\)--\(40\) (\(N=4\)
is the degenerate case \(L=1\), outside the scope of this mechanism =
3-wave closure by a different mechanism, §6).

\textbf{Why ``8'' (2-adic selection rule).} The original cyclic phase
advances by \(\Delta\varphi=4\pi/N\) under \(i\to i+1\), and the square
map doubles it again to \(\Delta\arg(z^2)=8\pi/N\). A two-wave
zero-closure \(w_b=-w_a\) requires an antipodal point (a subgroup of
order 2) in the cyclic group \(\mathbb Z_L\) of the square waves, which
holds iff \(L\) is even \(\iff\nu_2(N)\ge3\iff 8\mid N\) (\(\nu_2\) is
the 2-adic valuation). That is, \(N=8k\) is not an accidental period but
a \textbf{2-adic selection rule required by the square map}.

\textbf{Explicit construction for \(N=8k\) (no graph search needed).}
With \(L=N/4=2k\) and \(m=N/8\), we have \(q^m=e^{i\pi}=-1\). Hence
within the same class

\[w_{i+m}=-w_i\ \Longrightarrow\ z_{i+m}^2=-z_i^2\ \Longrightarrow\ z_{i+m}=\pm i\,z_i\]

(equal amplitude \(r_d\), phase difference exactly \(\pi/2\)). Pairing
the antipodal points of the map \(i\mapsto i+m\) (\(w_{i+m}=-w_i\)) in
each distance class produces a perfect matching on the spot (the orbit
length of \(i\mapsto i+m\) is 8 for an ordinary class \(d<N/2\) and 4
for the diameter class \(d=N/2\); both being even, the orbit can be
partitioned into pairs without surplus or deficit). For the \(8k\)
values in \(N=8\)--\(40\) we confirmed exact \(\pm i\) and a covering of
all \(M\) edges into \(M/2\) pairs (\(14,60,138,248,390\)). The perfect
matching of §4 has this \textbf{explicit construction}.

\textbf{Even including across distinct distance classes,
\(\iff 8\mid N\).} A \(\pm i\) pair requires \(|z_a|=|z_b|\) and a phase
difference of exactly \(\pi/2\). (1) Since the floor phases lie on
\((2\pi/N)\mathbb Z\), a \(\pi/2\) phase difference requires \(4\mid N\)
(impossible for \(4\nmid N\) = odd and \(N\equiv2\ (\mathrm{mod}\,4)\)).
(2) Since the \(2i\) in the phase \(2\pi(2i+d)/N\) is even, a \(\pi/2\)
difference between classes \(d,d'\) requires
\(d'-d\equiv N/4\ (\mathrm{mod}\,2)\). (3) Equal amplitude holds under
reflection \(r_d=r_{N/2-d}\) (numerically confirmed for \(N=3\)--\(40\),
error \(\le1.2\times10^{-16}\)), which gives
\(d'-d\equiv N/2\ (\mathrm{mod}\,2)\): for
\(N\equiv4\ (\mathrm{mod}\,8)\), \(N/2\) is even but \(N/4\) is odd, so
the two are incompatible, and only \(8\mid N\) is realized. We confirmed
that the count of exact \(\pm i\) pairs over all pairs (including across
classes) agrees perfectly and independently with the earlier exhaustive
audit (§4). Below, as consequences of this theorem, we describe the even
system (§4), the odd system (§5), and distance-class-independent closure
(§6).

\subsection{\texorpdfstring{4. Even Systems --- Complete \(\pm i\)
Photon-Pair Covering of
\(N=8k\)}{4. Even Systems --- Complete \textbackslash pm i Photon-Pair Covering of N=8k}}\label{even-systems-complete-pm-i-photon-pair-covering-of-n8k}

\begin{itemize}
\tightlist
\item
  \textbf{Exact \(z_b=\pm i z_a\) exists only for \(N=8k\)
  (8,16,24,32,40)} (other even \(N\) have 0 candidates)------a
  consequence of \(p=2\iff8\mid N\) in the theorem of §3. The following
  is its combinatorial structure.
\item
  All \(M\) waves can be \textbf{perfectly matched} into \(M/2\)
  non-overlapping two-wave self-closure pairs: \(N=8{:}14/28\),
  \(16{:}60/120\), \(24{:}138/276\), \(32{:}248/496\), \(40{:}390/780\).
\item
  Closed form of the candidate edge count
  \(E_\gamma=N(4N-11)/2=84,424,1020,1872,2980\). The candidate graph is
  \textbf{bipartite} for all \(N=8k\), with degrees only \(2/4/8\)
  (\(\#\deg2=N/2\), \(\#\deg4=N\), \(\#\deg8=M-3N/2\)) and
  connected-component sizes only \(4/8/16\) (\(C_4=N/8\), \(C_8=N/8\),
  \(C_{16}=N(N-4)/32\)).
\item
  \textbf{Selection rule (parity):} a \(90^\circ\) phase requires
  \(\Delta(i+j)=N/4\), and an edge of distance \(d\) satisfies
  \(i+j\equiv d\ (\mathrm{mod}\,2)\). If \(N=8k\), then \(N/4\) is even
  and parity is preserved; if \(N=8k+4\), it is odd and parity is
  inverted, incompatible with the equal-amplitude class.
\item
  \textbf{Complete pair decomposition is non-unique:} a perfect matching
  can select pairs sharing an original edge in anywhere from 0 up to a
  maximum of \(N\) pairs. From the static state alone, ``which two waves
  are one photon'' is not unique, and the absolute counts of \(+i/-i\)
  (right-handed/left-handed) cannot be defined from undirected pairs
  alone (propagation direction / orientation is separately required).
\end{itemize}

\subsection{\texorpdfstring{5. Odd Systems --- Minimal \(p\)-Wave
Condensate}{5. Odd Systems --- Minimal p-Wave Condensate}}\label{odd-systems-minimal-p-wave-condensate}

\begin{itemize}
\tightlist
\item
  For odd \(N\) (3,5,\ldots,39), \textbf{photon-type two-wave pairs are
  completely 0} (by the theorem of §3, odd cases have \(L=N\),
  \(p=\operatorname{spf}(N)\ge3\)), maximum matching 0, and all \(M\)
  waves are unpaired. Nevertheless, a \textbf{100\% complete
  decomposition} into minimal \(p\)-wave zero-closure condensates is
  possible (constructive coverage \(=1.0\), maximum group residual
  \(\le1.2\times10^{-15}\)).
\item
  \(z^2\) makes one round of the \(N\)-th roots and splits into regular
  \(p\)-gons with \(p=\text{smallest prime factor}(N)\):
  \(3\mid N\to3\)-wave, \(N=5,25,35\to\) minimal 5-wave, odd prime
  \(N\to\) minimal \(N\)-wave.
\item
  Over all 3-wave exhaustive checks, exact triads appear only for
  \(3\mid N\) (\(N=3,9,15,21,27,33,39\)). All 4-waves are 0 across all
  odd systems.
\end{itemize}

\subsection{6. Residual Waves and Independent Closure of Distance
Classes (Even
Systems)}\label{residual-waves-and-independent-closure-of-distance-classes-even-systems}

\begin{itemize}
\tightlist
\item
  \(N=8k\) has 0 residual (complete pair covering); for other even
  \(N\), exact 2-wave pairs are 0 and the residual = all \(M\) waves.
  \textbf{The residual as a whole is also zero-closed to machine
  precision} (relative residual
  \(4.1\times10^{-17}\)--\(2.4\times10^{-16}\)).
\item
  For \(N>4\), \textbf{each distance class \(C_d\) itself independently
  satisfies \(\sum_{e\in C_d}z_e^2=0\)}. The total closure is already
  decomposed into distance-class condensates. The intra-class phase step
  of the square wave \(w=z^2\) is \(q=e^{i8\pi/N}\), with order
  \(L=N/\gcd(N,4)\). For each factor \(p\) of \(L\), it can be
  decomposed into a \(p\)-gon zero-closure.
\item
  For all even \(N\), all \(M\) waves are 100\% decomposed into minimal
  symmetric zero-closure blocks with no overlap
  (e.g.~\(N=4{:}3\times2\), \(6{:}3\times5\), \(8{:}2\times14\),
  \(10{:}5\times9\), \(12{:}3\times22\), \ldots, \(40{:}2\times390\)).
  Exact 3-wave closure holds for \(N=4,6,12,18,24,30,36\). The
  factorization is non-unique when \(L\) is composite (\(N=24{:}2/3/6\),
  \(30{:}3/5/15\), \(40{:}2/5/10\)).
\end{itemize}

\subsection{7. On the Floor It Is a Single Carrier --- No
Harmonics}\label{on-the-floor-it-is-a-single-carrier-no-harmonics}

We audit the positive eigenvalues of \(iK\) of the frozen generator
\(K\) of the very cyclic-Fourier floor factorized in this paper (§2,
\(z_{ij}=r_d e^{i2\pi(i+j)/N}\)) (the audited floor is identical to
§2--6, reconstructed via \texttt{build\_theoretical\_floor} from
\(z=r_d e^{i\theta}\), \(\theta=2\pi(i+j)/N\); distinct from the
90-degree integer-K floor = Paper 3):

\begin{itemize}
\tightlist
\item
  The occupied state of the high-symmetry floor is a \textbf{single
  eigenmode} in which all waves co-wind a single carrier (eigen-residual
  \(\le3.0\times10^{-14}\)). Even after factorization, the inter-block
  frequency ratio on the floor is \(1{:}1\).
\item
  Across all even systems \(N=4\)--\(40\), there are \textbf{0
  integer-harmonic relations of 2/3/4-fold, 0 low-denominator rational
  ratios with denominator \(\le16\), and 0 three-wave sum resonances
  \(\omega_i+\omega_j=\omega_k\)}. The complete photon-pair covering of
  \(N=8k\) and the appearance of harmonics are uncorrelated.
\item
  A two-wave \(z_j=\pm iz_i\) is a zero-closure geometry with a
  \(90^\circ\) phase difference, not a harmonic such as \(1{:}2\). The
  high-symmetry background can be factorized into many closure
  condensates, but on the floor they have no independent clocks and are
  phase-locked to the same carrier.
\end{itemize}

\textbf{Control (what creates the factorization).} For the make\_parent
floor, which likewise satisfies self-consistency and zero-closure
\(\sum z^2=0\) and has a 90-degree phase lattice, we audited all
two-waves (Experiment 15), but exact \(z_b=\pm i\,z_a\) pairs are 0
across the entire series of \(N=3\)--\(40\) at \(\varepsilon\le10^{-9}\)
(\(\pm90^\circ\) phases form spontaneously, but amplitude equalization
is insufficient, with a best case of \(\varepsilon=3.7\times10^{-8}\)).
Therefore the photon-type two-wave factorization \textbf{does not arise
from zero-closure itself or from the 90-degree phase lattice alone}, but
depends on the phase-cyclicity (\(q=e^{i8\pi/N}\)) and the
distance-class equal-amplitude structure of the cyclic-Fourier floor
treated in this paper. This control is also consistent with the
make\_parent floor and the high-symmetry theoretical floor being
\textbf{distinct members} that satisfy the same ignition qualification
(Papers 3 and 0).

\subsection{8. Claim Classification and Open
Points}\label{claim-classification-and-open-points}

\begin{itemize}
\tightlist
\item
  Classification: \textbf{existence} of a static algebraic
  factorization. Within the same distance class it is a general proof
  from the phase alone (theorem of §3, \(L=N/\gcd(N,4)\),
  \(p=\operatorname{spf}(L)\)); the explicit construction for \(N=8k\)
  is exact; the full-space ``\(\iff 8\mid N\)'' is explained by the
  phase lattice + parity + amplitude reflection and numerically
  established by an exhaustive pair audit (\(N=3\)--\(40\)). Judgment:
  retained.
\item
  Open: that the equality of amplitudes is limited to the reflection
  \(r_d=r_{N/2-d}\) (that there is no accidental equal-amplitude class
  pair with odd \(d'-d\)) is numerically confirmed for \(N=3\)--\(40\).
  A general proof for arbitrary \(N\) (that the \(r_d\) of the
  cyclic-Fourier relative equilibrium has no equality other than the
  reflection) is left as a future task. However, the absence of 2-wave
  pairs for \(4\nmid N\) follows unconditionally from the phase lattice,
  independent of amplitude.
\item
  Not treated explicitly: identification of particle species and
  dynamical conservation (generation, annihilation, and recombination of
  blocks under fixed \(R_{124,23}\) exchange scattering) are outside the
  scope of this paper. The helicity of \(\pm i\) pairs (the absolute
  counts of right-handed/left-handed) cannot be defined from static
  undirected pairs alone.
\item
  Control: on the make\_parent floor, exact \(\pm i\) pairs \(=0\)
  (exhaustive pair audit, \(N=3\)--\(40\), Experiment 15). The
  factorization depends not on zero-closure or the 90-degree lattice but
  on the additional structure of the cyclic-Fourier floor (§7).
\item
  The high-symmetry background is more generally a \textbf{sea of
  minimal zero-closure condensates} than a ``sea of photons,'' and
  \(N=8k\) is the special series whose minimal unit can be decomposed
  down to two waves (photon-type). The algebraic basis for the minimal
  unit of information being a two-wave is Paper 9.
\end{itemize}

\subsection{Related Work (Structural Agreement, Not the Source of
Derivation)}\label{related-work-structural-agreement-not-the-source-of-derivation}

\begin{itemize}
\tightlist
\item
  Factorization of vanishing sums of roots of unity (Lam--Leung): agrees
  with the structure in which zero-closure blocks are decomposed into
  prime sizes.
\item
  Line graph / association scheme (\(L(K_N)\), Johnson scheme):
  distance-class structure.
\end{itemize}

\subsection{References}\label{references}

\begin{enumerate}
\def\labelenumi{\arabic{enumi}.}
\tightlist
\item
  Noriaki Kihara, ``The mechanism of inflationary rapid expansion in
  self-consistent relational-wave closed systems,'' Concept DOI
  10.5281/zenodo.22112008.
\item
  T. Y. Lam and K. H. Leung, ``On vanishing sums of roots of unity'', J.
  Algebra 224 (2000) 91.
\item
  D. Cvetković, P. Rowlinson, S. Simić, \emph{An Introduction to the
  Theory of Graph Spectra}, Cambridge (2010).
\end{enumerate}

\subsection{Reproduction}\label{reproduction}

Audit scripts and CSVs are provided as a complete set in each package:
for the even system,
\texttt{高対称理論床\_偶数系光子対精密解析\_20260908/}
(\texttt{audit\_highsym\_even\_photon\_pairs\_v2.py},
\texttt{even\_N\_summary.csv},
\texttt{selected\_maximum\_matchings.csv},
\texttt{component\_size\_counts.csv}),
\texttt{高対称理論床\_偶数系\_光子対と凝縮体因数分解\_20260908/}
(\texttt{analyze\_unpaired\_condensates\_v1.py},
\texttt{primitive\_zero\_closure\_groups.csv},
\texttt{distance\_class\_zero\_closure.csv}),
\texttt{高対称理論床\_偶数系\_倍音関係監査\_20260908/}
(\texttt{analyze\_highsym\_harmonic\_spectrum\_v2.py},
\texttt{positive\_frequency\_modes.csv}); for the odd system,
\texttt{高対称理論床\_奇数系\_ゼロ閉塞因数分解監査\_20260908/}
(\texttt{analyze\_highsym\_odd\_closure\_factorization\_v1.py},
\texttt{constructive\_primitive\_zero\_closure\_groups.csv}); and the
photon-pair audit of the make\_parent floor
\texttt{make\_parent型初期値\_自己無撞着構造\_20260907/二波自己閉塞光子対監査\_20260908/}.
The verification of the factorization theorem of §3 (\(L=N/\gcd(N,4)\),
\(p=\operatorname{spf}(L)\), the \(N=8k\) explicit construction, and
inter-class absence) is in \texttt{巡回二乗波因数分解定理\_20260913/}
(\texttt{verify\_cyclic\_square\_wave\_factorization\_20260913.py} =
read-out only, reproducing the phase theorem + the \(N=8k\) explicit
construction + the all-pairs audit,
\texttt{results/cyclic\_factorization\_summary.csv}, \texttt{実行ログ},
\texttt{分析}, README, SHA256SUMS, run\_all; the input is the
even-system CSVs above). Each is complete with README, SHA256SUMS,
analysis md, and RUN\_METADATA. The reconstruction and cross-checks of
the floor use the same definitions as the canonical
\texttt{run\_N3\_N40\_stage123\_v1.py} (SHA256
\texttt{1abf2353fee2e4f56f05e7a6f149fd086885136beb61ab571b48a56b09691567}),
with no physical changes.

\end{document}
